\documentclass[../DoAn.tex]{subfiles}
\begin{document}

\begin{center}
    \Large{\textbf{TÓM TẮT NỘI DUNG ĐỒ ÁN}}\\
\end{center}
\vspace{1cm}

Phát hiện xâm nhập mạng là bài toán phân loại lưu lượng mạng thành Benign và các loại tấn công trong thời gian thực. Hướng tiếp cận dựa trên dấu hiệu như Snort phát hiện hiệu quả các tấn công đã biết nhưng không phát hiện được biến thể mới và tấn công hành vi. Hướng tiếp cận dựa trên học máy cho phép học đặc trưng từ dữ liệu lịch sử, nhưng gặp hai thách thức chính trong triển khai thực tế: mất cân bằng lớp cực đoan (Benign chiếm trên 80\% lưu lượng) và covariate shift khi mô hình huấn luyện trên dữ liệu lab được triển khai sang môi trường mạng khác.

Đồ án này phát triển hệ thống phát hiện xâm nhập mạng lai ghép kết hợp Snort và học máy, được xây dựng qua ba giai đoạn nghiên cứu tuần tự. Giai đoạn 1 thiết lập kiến trúc nền tảng trên bộ dữ liệu NSL-KDD: Autoencoder Anomaly Gate lọc lưu lượng bình thường ở tầng thứ nhất, Stacking Ensemble gồm FT-Transformer kết hợp LightGBM qua Meta Logistic Regression phân loại tấn công ở tầng thứ hai. Mô hình đạt Macro F1 = 0,6809 trên KDDTest+, mức cạnh tranh với đánh giá trung thực theo giao thức train/test tách biệt. Giai đoạn 2 tái thiết kế kiến trúc Two-Stage Cascade cho bộ dữ liệu CIC-IDS-2017: Gating Network phân loại 6 lớp (Benign, DoS, DDoS, PortScan, Brute Force, Suspicious) với ngưỡng 85\% cho lớp Suspicious, Expert Network phân loại 5 lớp (Benign, Web Attack, Bot, Infiltration, Heartbleed) cho các flow Suspicious; hệ thống xuất tổng hợp 9 lớp. Asymmetric Ensemble Voting kết hợp FT-Transformer, Random Forest và KNN với luật bỏ phiếu bất đối xứng cho Botnet và Infiltration; Hard Negative Mining hai vòng cải thiện Botnet F1 từ 0,4787 lên 0,6512. Hệ thống đạt Accuracy 99,55\% và Macro F1 = 0,9294 trên tập kiểm tra CIC-IDS-2017. Giai đoạn 3 giải quyết vấn đề covariate shift khi triển khai mô hình sang môi trường Testbed WSL2: mô hình CIC giảm từ 99,55\% xuống 21,93\% Accuracy khi áp dụng trực tiếp. Sau khi thu thập dữ liệu tấn công thực từ hai máy trên LAN và tái huấn luyện hoàn toàn, mô hình V8.5 đạt Macro F1 = 91,7\% trên 5 lớp với BruteForce Recall = 90\% trên tập kiểm định đa dạng miền.

Đồ án đóng góp: (1) kiến trúc Two-Stage Cascade với Asymmetric Ensemble Voting cho bài toán phân loại đa lớp mất cân bằng cực đoan; (2) Hard Negative Mining hai vòng cải thiện F1 lớp thiểu số vượt trội class weight đơn thuần; (3) quy trình chẩn đoán covariate shift định lượng và tái huấn luyện theo dữ liệu đích; (4) bằng chứng thực nghiệm cho thấy tập kiểm định đồng nhất miền có thể inflate metric lên 14\% tuyệt đối so với tập kiểm định đa dạng miền.

\vspace{1.5cm}
\begin{flushright}
\textit{Sinh viên thực hiện}\\[1.5cm]
\textit{Đỗ Tuấn Minh}
\end{flushright}

\end{document}
